\chapter{Releases, Jailfiles, and Images}
\label{ch:images}

Blackship separates the base FreeBSD userland from application-specific image
construction. Releases are reusable roots; Jailfiles describe deterministic
changes; a running or built jail can be committed back into another release.

\section{Bootstrapping releases}

\begin{lstlisting}[language=bash]
blackship bootstrap 15.1-RELEASE
blackship bootstrap 15.1-RELEASE --archives base,lib32
blackship bootstrap 15.1-RELEASE --pkgbase
blackship releases list
blackship releases verify 15.1-RELEASE
\end{lstlisting}

FreeBSD 16 and newer use pkgbase automatically. For pkgbase, repository
metadata is checked through the FreeBSD fingerprint trust chain. With ZFS, a
successful bootstrap creates a release dataset and \command{@pristine}
snapshot. Blackship refuses to treat an incomplete root as a valid release.

\section{Line-oriented Jailfiles}

The default format deliberately resembles a Dockerfile while retaining
FreeBSD semantics:

\begin{lstlisting}[language=Jailfile]
FROM 15.1-RELEASE
LABEL service=web version=1.0
ARG NGINX_VERSION=1.30
ENV NGINX_VERSION=${NGINX_VERSION}
RUN pkg install -y nginx
RUN sysrc nginx_enable=YES
COPY nginx.conf /usr/local/etc/nginx/nginx.conf
COPY public/ /usr/local/www/nginx/
WORKDIR /usr/local/www/nginx
EXPOSE 80/tcp
CMD /usr/local/sbin/nginx -g 'daemon off;'
STOP /usr/local/sbin/nginx -s quit
\end{lstlisting}

\begin{shiptable}{L{0.17\textwidth}L{0.73\textwidth}}
\toprule
Instruction & Meaning \\
\midrule
\command{FROM} & Required base release or committed image. \\
\command{LABEL} & Image metadata. \\
\command{ARG} & Build-time variable with an optional default. \\
\command{ENV} & Environment metadata available to later instructions. \\
\command{RUN} & Command executed while building the jail root. \\
\command{COPY} & Copy from the bounded build context into the jail root. \\
\command{WORKDIR} & Set/create the working directory for later steps. \\
\command{EXPOSE} & Document an intended service port; it does not create PF rules. \\
\command{CMD} & Default image command metadata. \\
\command{STOP} & Command run before rc shutdown; manifest \setting{stop} wins. \\
\bottomrule
\end{shiptable}

Source resolution refuses absolute paths and parent traversal. Destination
writes refuse symlink escapes from the jail root. Ownership is resolved against
the jail, not the host.

\section{TOML Jailfiles}

The alternative format is convenient for generated configuration:

\begin{lstlisting}[language=TOML]
[metadata]
name = "web"
version = "1.0"

from = "15.1-RELEASE"

[[args]]
name = "NGINX_VERSION"
default = "1.30"

[[instructions]]
type = "run"
command = "pkg install -y nginx"

[[instructions]]
type = "copy"
src = "nginx.conf"
dest = "/usr/local/etc/nginx/nginx.conf"

[[expose]]
port = 80
protocol = "tcp"

cmd = "/usr/local/sbin/nginx -g 'daemon off;'"

[stop]
cmd = "/usr/local/sbin/nginx -s quit"
\end{lstlisting}

Generate either form with \command{blackship init}; pass \command{--toml} for
the structured variant.

\section{Build, inspect, and validate}

\begin{lstlisting}[language=bash]
blackship template validate Jailfile
blackship template inspect Jailfile
blackship build --file Jailfile --name web --dry-run
blackship build --file Jailfile --name web \
  --build-arg NGINX_VERSION=1.30
\end{lstlisting}

Dry run does not bootstrap a missing release or mutate storage. Validation
should be part of CI before a privileged build job.

\section{ZFS layer cache}

Every cacheable \command{RUN} and \command{COPY} contributes to a rolling hash
of the base release, earlier instructions, build arguments, and copied content.
A completed step is snapshotted as \command{bs-layer-N-hash}. A rebuild finds
the longest matching prefix, rolls back to it, and resumes at the first changed
step.

Changing an early instruction invalidates it and all later layers. Changing an
unrelated file outside a \command{COPY} source does not. This makes build
contexts both faster and more predictable than timestamp-only caching.

\section{Commit and reuse}

\begin{lstlisting}[language=bash]
blackship build --file runner/Jailfile --name runner
blackship commit runner runner-golden
blackship run --name ci-42 --release runner-golden -- ./run-tests.sh
\end{lstlisting}

Commit stages the known jail definition and Jailfile under
\file{.blackship/}, snapshots the root, clones it into the release tree, and
promotes it so the result is independently reusable. Choose immutable,
meaningful release names; consumers should not need to know whether the seed
came from FreeBSD or a committed image.

\section{Armada build references}

\begin{lstlisting}[language=TOML]
[[jails]]
name = "web"
build = "./web"
depends_on = ["database"]

[[jails]]
name = "worker"
jailfile = "./worker/Worker.jailfile"
\end{lstlisting}

\setting{build} names a directory containing \file{Jailfile};
\setting{jailfile} names an explicit file. Armada's \command{up} builds missing
roots by default, \command{--build} forces a rebuild, and
\command{--no-build} fails when a required root does not already exist.

\begin{shipnote}[Builds are not runtime configuration]
Put packages, files, and service enablement in a Jailfile. Put host attachment,
resource policy, addresses, dependencies, and operational hooks in the
manifest. Keeping the boundary clear makes images reusable and manifests
auditable.
\end{shipnote}
